% =====================================================================
% Background — Walsh--Hadamard transform, coalition sweeps, discovery
% methods, MIB faithfulness metrics.
%
% Status: v1 first draft. Compiles with main.tex.
% =====================================================================

\section{Background}
\label{sec:background}

\subsection{Walsh--Hadamard interaction decomposition}
\label{sec:walsh-decomposition}

Given a circuit of $n$ components (attention heads), the coalition
function $v: \{0,1\}^n \to \mathbb{R}$ maps each subset of active heads
to the model's performance metric (logit difference) when those heads
operate normally and the rest are ablated. The $2^n$ function values
admit a unique Walsh--Hadamard decomposition
\begin{equation}
  v(x) = \sum_{S \subseteq [n]} w_S \,\chi_S(x),
  \qquad
  \chi_S(x) = (-1)^{\sum_{i \in S} x_i},
\end{equation}
where the sum ranges over all $2^n$ subsets $S$ of $[n] = \{1, \ldots, n\}$
and $\chi_S$ is the corresponding Walsh basis
function~\citep{poelwijk2016context}. The coefficient $w_S$ is the
unique order-$|S|$ interaction among the heads in $S$:
\begin{equation}
  w_S = \frac{1}{2^n} \sum_{x \in \{0,1\}^n} v(x)\,\chi_S(x).
  \label{eq:walsh-coeff}
\end{equation}
Order-1 coefficients $w_{\{i\}}$ measure each head's marginal
importance averaged over all $2^{n-1}$ contexts. Order-2 coefficients
$w_{\{i,j\}}$ measure how much the joint effect of heads $i$ and $j$
deviates from the sum of their individual effects, again averaged over
all contexts. The decomposition is reference-free: unlike pairwise
epistasis measured against a single wildtype, each Walsh coefficient
averages the relevant interaction over all possible backgrounds,
eliminating dependence on an arbitrary reference
state~\citep{poelwijk2019learning}.

The transform satisfies Parseval's identity: $\sum_S w_S^2 = 2^{-n}
\sum_x v(x)^2$. The fraction of total variance (``energy'') carried by
each order quantifies how much of the circuit's behavior requires
combinations of that many heads. A circuit whose energy concentrates at
order 1 is well approximated by an additive model; one with substantial
order-2 energy has pairwise interactions that an additive summary
discards.

\subsection{Ablation primitives}
\label{sec:ablation-primitives}

The coalition function depends on how inactive heads are ablated.
Three primitives are standard:

\begin{itemize}
\item \textbf{Zero ablation}: replace the head's output
  (\texttt{hook\_z}) with zero. Removes the head's contribution to the
  residual stream entirely.

\item \textbf{Mean ablation}: replace \texttt{hook\_z} with its mean
  across prompts. Substitutes the head's average behavior, preserving
  the residual stream's scale.

\item \textbf{Resample ablation}: replace \texttt{hook\_z} with its
  value on a different prompt from the same distribution. Preserves
  distributional statistics while destroying task-specific information.
\end{itemize}

Different primitives yield quantitatively different Walsh spectra. We
report results under all three where exhaustive sweeps are available
(15-head circuits) and under mean ablation for the 20-head circuit.

\subsection{Circuit discovery methods}
\label{sec:discovery-methods}

\paragraph{Activation patching.}

The marginal importance of each head is measured by ablating it in
isolation and recording the change in the target metric. Activation
patching by mean ablation on \texttt{hook\_z} is the standard
protocol~\citep{wang2022ioi}. The resulting scores are order-1
quantities: they measure each head independently, discarding all
interaction structure.

\paragraph{Edge Attribution Patching (EAP).}

\citet{syed2023attribution} approximate the effect of ablating each
edge (sender--receiver pair) by backpropagating through the clean
forward pass and multiplying gradients by activation differences. The
resulting edge scores are linearized interaction estimates: they
capture how much one head's output affects another's computation,
limited to first order. EAP runs in a single forward--backward pass and
scales to large circuits.

\paragraph{Path patching.}

Path patching isolates the direct causal edge from sender $S$ to
receiver $R$ by corrupting $S$, freezing all intermediate components to
their clean values, and letting $R$ compute naturally from the modified
residual stream~\citep{wang2022ioi}. The change in the target metric
relative to the clean run is the direct effect of $S$ on the output
routed specifically through $R$. Path patching distinguishes direct
edges from mediated interactions but requires one forward pass per edge.

\subsection{MIB faithfulness metrics}
\label{sec:mib-metrics}

The Mechanistic Interpretability Benchmark
(MIB)~\citep{mueller2025mib} evaluates circuit discovery methods by how
faithfully a discovered circuit reproduces the full model's behavior.
Given a metric $m$ (such as logit difference), the faithfulness of a
circuit $C$ relative to the full model $N$ is
\begin{equation}
  f(C, N; m) = \frac{m(C) - m(\varnothing)}{m(N) - m(\varnothing)},
  \label{eq:faithfulness}
\end{equation}
where $m(\varnothing)$ is the metric with all components ablated and
$m(N)$ is the metric with no ablation. A faithfulness of 1 means the
circuit fully reproduces the model's behavior; 0 means the circuit
performs no better than complete ablation.

Two summary statistics aggregate the faithfulness curve across circuit
sizes:

\begin{itemize}
\item \textbf{CPR} (Circuit Performance Recovered): the area under the
  faithfulness curve as a function of circuit size, normalized to $[0,1]$.
  Higher is better.

\item \textbf{CMD} (Circuit--Model Distance): the area between the
  faithfulness curve and the ideal curve $f = 1$, normalized to $[0,1]$.
  Lower is better. CMD $= 1 - \text{CPR}$ when faithfulness is monotone.
\end{itemize}
